\documentclass[letterpaper,12pt,oneside]{article}
\usepackage[top=1in,bottom=1in,left=1.25in,right=1.25in]{geometry}
\usepackage[T1]{fontenc}
\usepackage[utf8]{inputenc}
\usepackage[spanish,es-nodecimaldot,es-tabla]{babel}
\usepackage{lmodern,microtype,graphicx,caption,pdfpages}
\usepackage{amsmath,amssymb,booktabs,tabularx,longtable,xcolor,adjustbox}
\usepackage[hidelinks]{hyperref}
\numberwithin{figure}{section}

\usepackage{needspace}
\usepackage[section]{placeins}
\setlength{\parindent}{0pt}
\setlength{\parskip}{0.45em}
\setlength{\emergencystretch}{2em}
\widowpenalty=10000
\clubpenalty=10000
\displaywidowpenalty=10000
\renewcommand{\topfraction}{0.95}
\renewcommand{\bottomfraction}{0.95}
\renewcommand{\textfraction}{0.05}
\renewcommand{\floatpagefraction}{0.80}
\renewcommand{\refname}{Bibliograf\'ia}

\ifdefined\PortadaPDF
\else
  \IfFileExists{portada_config.tex}{%
    \input{portada_config.tex}%
  }{}%
\fi

\newcommand{\figura}[3]{%
  \begin{figure}[!htbp]
    \centering
    \adjustimage{max width=0.82\textwidth,max totalheight=0.25\textheight,keepaspectratio}{#1}%
    \caption{#2}%
    \label{#3}%
  \end{figure}%
}

\title{\textbf{Reporte de Pr\'actica 8}\\\large Huellas de Carbono e Interfaces en Java}
\author{\textbf{Equipo 2} --- Grupo 7\\Facultad de Ingenier\'ia, Universidad Nacional Aut\'onoma de M\'exico}
\date{05 de noviembre de 2026}

\begin{document}

\ifdefined\PortadaPDF
  \includepdf[pages=-,fitpaper=false,pagecommand={}]{\PortadaPDF}
\else
  \PackageError{reporte}{Falta definir PortadaPDF}{Inyecte la ruta privada de la portada convertida al compilar mediante -Z search-path con portada_config.tex o definiendo \string\PortadaPDF.}
\fi

\tableofcontents
\clearpage

\section{Introducci\'on}

\subsection{Planteamiento del problema}
La pr\'actica solicita resolver dos necesidades fundamentales relacionadas con el paradigma orientado a objetos en Java. La primera consiste en plantear una interfaz gr\'afica interactiva para capturar un nombre de usuario y una contrase\~na mediante componentes de Swing. La segunda requiere modelar tres entidades diferenciadas ---un edificio, un autom\'ovil y una bicicleta--- a trav\'es de clases que implementen una misma interfaz denominada \texttt{HuellaCarbono}. Cada objeto debe estimar su huella a partir de variables propias de su dominio operativo, mientras que la clase principal debe organizar las instancias en una colecci\'on de referencias comunes para evidenciar el despacho polim\'orfico en tiempo de ejecuci\'on.

\subsection{Motivaci\'on}
El ejercicio permite contrastar el desacoplamiento que brindan las interfaces frente a la herencia tradicional. Una interfaz define un contrato de operaciones sin imponer estado compartido ni restringir la mec\'anica de c\'alculo a una f\'ormula \'unica. As\'i, el edificio vincula consumo el\'ectrico y factor de emisi\'on; el autom\'ovil procesa distancia recorrida, rendimiento de combustible y factor de gasolina; y la bicicleta expresa que su operaci\'on no produce emisiones directas en el alcance did\'actico delimitado. Por otra parte, la implementaci\'on de la interfaz gr\'afica complementa el aprendizaje al demostrar la manipulaci\'on segura de contrase\~nas y el despacho de eventos mediante oyentes.

\subsection{Objetivos}
\begin{itemize}
  \item Dise\~nar e implementar el contrato de la interfaz \texttt{HuellaCarbono} junto con sus clases concretas (\texttt{Edificio}, \texttt{Auto} y \texttt{Bicicleta}) para calcular emisiones seg\'un sus propios atributos y f\'ormulas de dominio.
  \item Demostrar el polimorfismo din\'amico en Java mediante el almacenamiento de instancias en una colecci\'on homog\'enea \texttt{ArrayList<HuellaCarbono>} y su invocaci\'on uniforme durante el recorrido.
  \item Construir una interfaz gr\'afica interactiva de inicio de sesi\'on con componentes Swing (\texttt{LoginGUI}) y verificar exhaustivamente la soluci\'on mediante compilaci\'on limpia, ejecuci\'on reproducible, generaci\'on de Javadoc y empaquetado en archivo JAR.
\end{itemize}

\section{Marco Te\'orico}

Una interfaz en Java constituye un contrato abstracto de operaciones que las clases concretas se comprometen a satisfacer. La especificaci\'on oficial del lenguaje establece que una clase puede implementar m\'ultiples interfaces y que cualquier referencia tipada a la interfaz puede emplearse para invocar sus m\'etodos declarados \cite{jls}. En esta pr\'actica, \texttt{HuellaCarbono} define \'unicamente la firma \texttt{double getHuellaCarbono()}, delegando enteramente en cada clase la l\'ogica de evaluaci\'on. La resoluci\'on del m\'etodo efectivo tiene lugar de manera din\'amica durante la ejecuci\'on, materializando el principio de ligadura tard\'ia y polimorfismo.

\Needspace{6\baselineskip}
El principio de encapsulamiento garantiza que el estado de cada entidad permanezca resguardado dentro de la clase correspondiente. De este modo, variables tales como consumo anual, kilometraje, rendimiento y factores de emisi\'on se declaran como atributos privados, proveyendo m\'etodos accesores p\'ublicos y la operaci\'on de c\'alculo exigida por el contrato. Se descarta deliberadamente el uso de herencia simple porque un edificio, un autom\'ovil y una bicicleta no comparten un antepasado sem\'antico com\'un en el mundo real; comparten exclusivamente una capacidad operativa de medici\'on ambiental.

En el \'ambito gr\'afico, la biblioteca Swing ofrece un conjunto integral de componentes de interfaz. La ventana principal se construye mediante un contenedor de primer nivel \texttt{JFrame}; la entrada de texto plano se gestiona con \texttt{JTextField}; el campo especializado \texttt{JPasswordField} garantiza el tratamiento confidencial de credenciales; el bot\'on \texttt{JButton} desencadena la acci\'on del usuario; y el componente \texttt{JLabel} retroalimenta visualmente el estado del acceso. La documentaci\'on de referencia de Oracle enfatiza la vinculaci\'on de controles con receptores de eventos a trav\'es de la interfaz \texttt{ActionListener} \cite{swing}.

La estimaci\'on de emisiones de gases de efecto invernadero depende del alcance metodol\'ogico y de los factores aplicables. Los inventarios internacionales categorizan las fuentes seg\'un emisiones directas e indirectas, donde el factor el\'ectrico responde a la mezcla de generaci\'on regional y temporal \cite{epaelectric}. A fin de conservar la reproducibilidad y verificabilidad del ejercicio, se definen dos factores did\'acticos estables: $0.40\ \mathrm{kg\ CO_2e/kWh}$ para suministro el\'ectrico y $2.31\ \mathrm{kg\ CO_2e/L}$ para combustible regular, valores derivados y redondeados a partir de inventarios t\'ecnicos oficiales de emisiones \cite{epafactors}.

\clearpage
\section{Desarrollo}

\subsection{Investigaci\'on y an\'alisis de la consulta}
El desarrollo de la pr\'actica tom\'o como punto de partida el an\'alisis formal de los requerimientos descritos en la especificaci\'on t\'ecnica (\texttt{doc/Capturas/query.txt}). Dicha formulaci\'on delimita la problem\'atica a resolver: dise\~nar un contrato \'unico \texttt{HuellaCarbono} capaz de desacoplar el c\'alculo num\'erico de emisiones entre tipos heterog\'eneos, y construir un m\'odulo gr\'afico interactivo de autenticaci\'on con componentes est\'andar de Swing.

La especificaci\'on formal fija los par\'ametros can\'onicos de prueba:
\begin{itemize}
  \item \textbf{Edificio:} Evaluado con un consumo el\'ectrico anual de $2400.00\ \mathrm{kWh}$ y un factor de conversi\'on regional de $0.40\ \mathrm{kg\ CO_2e/kWh}$, produciendo una emisi\'on anal\'itica de $960.00\ \mathrm{kg\ CO_2e}$.
  \item \textbf{Autom\'ovil:} Evaluado sobre un trayecto de $480.00\ \mathrm{km}$, un consumo nominal de $7.00\ \mathrm{L/100\,km}$ y un factor de combusti\'on de gasolina de $2.31\ \mathrm{kg\ CO_2e/L}$. El consumo neto corresponde a $33.60\ \mathrm{L}$, generando una huella anal\'itica de $77.616\ \mathrm{kg\ CO_2e}$ (redondeada reglamentariamente a $77.62\ \mathrm{kg\ CO_2e}$).
  \item \textbf{Bicicleta:} Modelada sobre un recorrido de $12.00\ \mathrm{km}$ bajo el alcance delimitado de emisiones directas operativas nulas ($0.00\ \mathrm{kg\ CO_2e}$).
  \item \textbf{Caso frontera:} Incorporaci\'on voluntaria de un edificio con consumo nulo ($0.00\ \mathrm{kWh}$) para auditar la estabilidad aritm\'etica de la soluci\'on ante condiciones de borde.
\end{itemize}

El dise\~no del sistema garantiza la cohesi\'on y separaci\'on de responsabilidades, organizando las entidades bajo el paquete institucional \texttt{mx.unam.fi.die.poo.g7.p8}.

\clearpage
\subsection{Interfaz gr\'afica de login}
La clase \texttt{LoginGUI} extiende de \texttt{JFrame} e implementa de manera directa la interfaz \texttt{ActionListener}. El formulario organiza un esquema matricial compuesto por un campo de texto para usuario, un campo protegido para contrase\~na (\texttt{JPasswordField}), un bot\'on de ingreso (\texttt{JButton}) y una etiqueta de estado (\texttt{JLabel}). La arquitectura est\'atica de la clase se presenta en el diagrama de clases de la Figura~\ref{fig:loginuml}, mientras que la din\'amica de interacci\'on y captura de eventos se ilustra en el diagrama de flujo de la Figura~\ref{fig:loginflujo}. La validaci\'on compara las entradas frente a las credenciales did\'acticas (\texttt{usuario} y clave \texttt{poo8}), notificando el veredicto en la interfaz.

\figura{doc/Diagramas/LoginGUI_uml.png}{Diagrama UML de la clase LoginGUI.}{fig:loginuml}
\figura{doc/Diagramas/LoginGUI_flujo.png}{Flujo de captura, eventos y validaci\'on del inicio de sesi\'on.}{fig:loginflujo}
\FloatBarrier

\clearpage
\subsection{Contrato de huella de carbono}
La interfaz \texttt{HuellaCarbono} estipula un contrato consistente en un \'unico m\'etodo: \texttt{double getHuellaCarbono()}. Al mantener la interfaz concisa y altamente cohesiva, los componentes clientes interact\'uan exclusivamente a trav\'es de la abstracci\'on sin acoplarse a las estructuras internas de cada implementaci\'on. Como se aprecia en la Figura~\ref{fig:huellauml}, el compartimento de atributos permanece vac\'io debido a que la interfaz no declara estado mutable. Por su parte, la Figura~\ref{fig:huellaflujo} modela el flujo operacional por el cual cada tipo concreto asume la obligaci\'on de resolver el c\'alculo.

\figura{doc/Diagramas/HuellaCarbono_uml.png}{Diagrama UML de la interfaz HuellaCarbono.}{fig:huellauml}
\figura{doc/Diagramas/HuellaCarbono_flujo.png}{Flujo conceptual de delegaci\'on de la interfaz HuellaCarbono.}{fig:huellaflujo}
\FloatBarrier

\clearpage
\subsection{Clase Edificio}
La clase \texttt{Edificio} cuantifica el impacto ambiental a partir del consumo anual de energ\'ia el\'ectrica y su factor de emisi\'on correspondiente. La relaci\'on de c\'alculo adoptada es directa: $H = E \times F_e$, donde $E$ denota los kilowatt-hora acumulados y $F_e$ el factor de emisi\'on el\'ectrica. Los datos miembros se mantienen bajo protecci\'on privada (\texttt{nombre}, \texttt{consumoAnualKwh} y \texttt{factorElectricidad}), proveyendo m\'etodos accesores expl\'icitos conforme al diagrama UML de la Figura~\ref{fig:edificiouml}. El algoritmo de evaluaci\'on de la huella se detalla en el diagrama de flujo de la Figura~\ref{fig:edificioflujo}.

\figura{doc/Diagramas/Edificio_uml.png}{Diagrama UML de la clase Edificio.}{fig:edificiouml}
\figura{doc/Diagramas/Edificio_flujo.png}{Flujo de c\'alculo de huella el\'ectrica en la clase Edificio.}{fig:edificioflujo}
\FloatBarrier

\clearpage
\subsection{Clase Auto}
La clase \texttt{Auto} representa un veh\'iculo terrestre automotor. El volumen de combustible quemado se eval\'ua preliminarmente mediante $L = (K \times R) / 100$, donde $K$ denota la distancia en kil\'ometros y $R$ el rendimiento expresado en litros por cada cien kil\'ometros. Con dicho valor, la huella total se obtiene con la relaci\'on $H = L \times F_g$, empleando el factor de combusti\'on de gasolina $F_g$. En la Figura~\ref{fig:autouml} se observa el encapsulamiento de sus atributos y en la Figura~\ref{fig:autoflujo} se diagrama la secuencia aritm\'etica de dos fases que favorece la trazabilidad del c\'alculo.

\figura{doc/Diagramas/Auto_uml.png}{Diagrama UML de la clase Auto.}{fig:autouml}
\figura{doc/Diagramas/Auto_flujo.png}{Flujo de c\'alculo de combustible y huella de la clase Auto.}{fig:autoflujo}
\FloatBarrier

\clearpage
\subsection{Clase Bicicleta}
La clase \texttt{Bicicleta} gestiona el tipo de unidad y el recorrido acumulado. En el modelo educativo formulado, el alcance se restringe exclusivamente a las emisiones operativas directas generadas durante el trayecto, las cuales resultan nulas al no existir combusti\'on interna ($H = K \times 0.00 = 0$). La estructura de atributos encapsulados se exhibe en la Figura~\ref{fig:bicicletouml}, y el diagrama de flujo correspondiente a su evaluaci\'on directa se ilustra en la Figura~\ref{fig:bicicletaflujo}.

\figura{doc/Diagramas/Bicicleta_uml.png}{Diagrama UML de la clase Bicicleta.}{fig:bicicletouml}
\figura{doc/Diagramas/Bicicleta_flujo.png}{Flujo de evaluaci\'on de emisiones operativas de Bicicleta.}{fig:bicicletaflujo}
\FloatBarrier

\clearpage
\subsection{Clase principal y despacho polim\'orfico}
La clase \texttt{Main} materializa el escenario de prueba instanciando los tres elementos solicitados, adem\'as de un caso de prueba frontera consistente en un edificio con consumo el\'ectrico nulo. Las cuatro instancias se agregan a una lista gen\'erica de tipo \texttt{ArrayList<HuellaCarbono>}. Al recorrer la colecci\'on, el programa interact\'ua \'unicamente con el tipo general de la interfaz: invoca \texttt{getHuellaCarbono()} y la m\'aquina virtual resuelve internamente el m\'etodo especializado en tiempo de ejecuci\'on. La Figura~\ref{fig:mainuml} presenta la composici\'on est\'atica de la clase principal y la Figura~\ref{fig:mainflujo} esquematiza el ciclo de recorrido polim\'orfico.

\figura{doc/Diagramas/Main_uml.png}{Diagrama UML de la clase principal Main.}{fig:mainuml}
\figura{doc/Diagramas/Main_flujo.png}{Flujo de instanciaci\'on, colecci\'on y recorrido polim\'orfico.}{fig:mainflujo}
\FloatBarrier

\clearpage
\section{Resultados}

\subsection{Verificaci\'on cuantitativa de emisiones}
Las pruebas experimentales se contrastaron frente a los c\'alculos anal\'iticos te\'oricos previamente formulados. El edificio con consumo anual de $2400.00\ \mathrm{kWh}$ y factor $0.40\ \mathrm{kg\ CO_2e/kWh}$ reporta exactamente $960.00\ \mathrm{kg\ CO_2e}$. El autom\'ovil que recorri\'o $480.00\ \mathrm{km}$ a $7.00\ \mathrm{L/100\,km}$ consumi\'o $33.60\ \mathrm{L}$ de gasolina; multiplicado por el factor $2.31\ \mathrm{kg\ CO_2e/L}$ arroj\'o $77.616\ \mathrm{kg\ CO_2e}$, registrando $77.62$ tras redondeo can\'onico a dos posiciones decimales. La bicicleta produce $0.00\ \mathrm{kg\ CO_2e}$ por definici\'on operativa directa, y el caso frontera con consumo cero arroja $0.00\ \mathrm{kg\ CO_2e}$, ratificando que consumos nulos no introducen valores espurios. Los resultados cuantitativos se compendian en el Cuadro~\ref{tab:resultados}.

\begin{table}[!htbp]
\centering
\caption{Comparaci\'on cuantitativa entre valores esperados y observados en ejecuci\'on.}
\label{tab:resultados}
\begin{adjustbox}{max width=\textwidth}
\begin{tabular}{p{0.24\textwidth}p{0.36\textwidth}rrp{0.14\textwidth}}
\toprule
\textbf{Instancia} & \textbf{Par\'ametros de entrada} & \textbf{Esperado} & \textbf{Observado} & \textbf{Dictamen} \\
\midrule
Edificio Escuela & 2400.00~kWh, factor $0.40$ & 960.00 & 960.00 & Exitoso \\
Auto Sed\'an & 480.00~km, 7.0~L/100\,km, factor $2.31$ & 77.62 & 77.62 & Exitoso \\
Bicicleta urbana & 12.00~km, factor operativo $0.00$ & 0.00 & 0.00 & Exitoso \\
Edificio frontera & 0.00~kWh, factor $0.40$ & 0.00 & 0.00 & Exitoso \\
\bottomrule
\end{tabular}
\end{adjustbox}
\end{table}

La ejecuci\'on comprob\'o la correcta emisi\'on textual de cada instancia y la resoluci\'on unificada de huella a trav\'es de la interfaz. La evidencia gr\'afica de la compilaci\'on limpia con el compilador \texttt{javac} se documenta en la Figura~\ref{fig:compilacion}, mientras que la verificaci\'on de empaquetado del archivo ejecutable \texttt{Practica8.jar} se muestra en la Figura~\ref{fig:jarcontenido}.

\figura{doc/Capturas/compilacion_javac.png}{Compilaci\'on limpia de los seis archivos fuente Java hacia la carpeta test.}{fig:compilacion}
\figura{doc/Capturas/empaquetado_jar.png}{Generaci\'on y verificaci\'on de contenido del archivo ejecutable Practica8.jar.}{fig:jarcontenido}

\Needspace{8\baselineskip}
\subsection{Ejecuci\'on desde classpath y archivo JAR}
La ejecuci\'on del programa se valid\'o bajo dos modalidades id\'enticas: directamente desde el \'arbol de clases compiladas en \texttt{test/} (Figura~\ref{fig:ejecucioncp}) y a trav\'es del archivo empaquetado \texttt{lib/Practica8.jar} provisto de manifiesto de ejecuci\'on (Figura~\ref{fig:ejecucionjar}). Ambas v\'ias arrojaron exactamente las cuatro lecturas evaluadas: $960.00$, $77.62$, $0.00$ y $0.00\ \mathrm{kg\ CO_2e}$, confirmando la paridad operativa entre ambos esquemas.

\figura{doc/Capturas/ejecucion_classpath.png}{Salida de consola observada al ejecutar Main desde la carpeta de clases compiladas.}{fig:ejecucioncp}
\figura{doc/Capturas/ejecucion_jar.png}{Salida de consola observada al ejecutar la aplicaci\'on desde Practica8.jar.}{fig:ejecucionjar}
\FloatBarrier

\subsection{Documentaci\'on t\'ecnica y estructura de entrega}
Se comprob\'o la generaci\'on independiente de documentaci\'on Javadoc para cada una de las seis entidades (\texttt{HuellaCarbono}, \texttt{Edificio}, \texttt{Auto}, \texttt{Bicicleta}, \texttt{Main} y \texttt{LoginGUI}) dentro de subdirectorios aislados en \texttt{doc/Javadoc}. La carpeta \texttt{test/} aloja exclusivamente los archivos binarios \texttt{.class} generados durante la compilaci\'on, y \texttt{lib/} conserva el archivo JAR aut\'onomo de la pr\'actica.

\clearpage
\section{Conclusiones}

La implementaci\'on de la interfaz \texttt{HuellaCarbono} confirm\'o la eficacia del polimorfismo por interfaz en Java: tres clases con modelos de dominio divergentes se articularon arm\'onicamente bajo un contrato unificado sin requerir jerarqu\'ias de herencia forzadas. Cada entidad resguard\'o sus atributos mediante encapsulamiento estricto y aplic\'o su propia f\'ormula de emisiones, permitiendo que la colecci\'on cliente opere de manera desacoplada y gen\'erica.

Los resultados experimentales observados en ejecuci\'on coincidieron plenamente con los valores anal\'iticos formulados antes de la prueba. El caso de frontera ratific\'o la estabilidad aritm\'etica de la implementaci\'on ante entradas nulas, mientras que las pruebas cruzadas entre classpath y archivo JAR garantizaron la reproducibilidad del entorno sin inconsistencias de configuraci\'on.

La clase \texttt{LoginGUI} complet\'o los requerimientos de interfaz gr\'afica interactiva demostrando el uso de campos protegidos \texttt{JPasswordField} y la delegaci\'on de eventos mediante \texttt{ActionListener}. En conjunto, la soluci\'on articula con precisi\'on los principios fundamentales de abstracci\'on, polimorfismo, verificaci\'on num\'erica y presentaci\'on formal de software de acuerdo con los criterios acad\'emicos de la asignatura.

\clearpage
\renewcommand{\refname}{Bibliograf\'ia}
\begin{thebibliography}{99}

\bibitem{jls}
Oracle Corporation, \emph{The Java Language Specification, Java SE 25 Edition: Interfaces}. Austin, TX, EE.~UU.: Oracle Corp., 2025.

\bibitem{swing}
Oracle Corporation, ``Creating a GUI With Swing'', en \emph{The Java Tutorials}. Redwood City, CA, EE.~UU.: Oracle Corp., 2024. [En l\'inea]. Disponible en: \url{https://docs.oracle.com/javase/tutorial/uiswing/}. [Accedido: 29-oct-2026].

\bibitem{epaelectric}
Agencia de Protecci\'on Ambiental de los Estados Unidos (EPA), \emph{Emissions \& Generation Resource Integrated Database (eGRID)}. Washington, D.C., EE.~UU.: EPA, 2024.

\bibitem{epafactors}
Agencia de Protecci\'on Ambiental de los Estados Unidos (EPA), \emph{Emission Factors for Greenhouse Gas Inventories}. Washington, D.C., EE.~UU.: EPA, 2024.

\end{thebibliography}

\end{document}
